\documentclass[a4,11pt]{aleph-notas}
% Se puede ver la documentación aquí: 
% https://github.com/alephsub0/LaTeX_aleph-notas 

% -- Paquetes adicionales 
\usepackage{enumitem}
\usepackage{array,booktabs,multirow,makecell}
\usepackage{colortbl}
\usepackage{longtable}
\usepackage{url}
\usepackage{pdflscape}
\usepackage{amsmath,amssymb}

% -- Datos 
\institucion{Facultad de Ciencias Exactas, Naturales y Ambientales}
\carrera{Catálogo STEM}
\asignatura{Álgebra Lineal y Geometría Analítica}
\tema{Cronograma y actividades}
\autor{Andrés Merino}
\fecha{Periodo 2026-1}

\logouno[0.16\textwidth]{Logos/logoPUCE_04_ac}
\definecolor{colortext}{HTML}{1957d1}
\definecolor{colordef}{HTML}{1957d1}
\fuente{montserrat}


% -- Comandos para tablas
\newcolumntype{C}[1]{>{\hspace{0pt}\centering\arraybackslash}p{#1}}
\newcolumntype{L}[1]{>{\raggedright\arraybackslash}p{#1}}

\definecolor{verde}{RGB}{0, 255, 127}
\definecolor{celeste}{RGB}{68,195,218}

\begin{document}
\addtolength{\headheight}{1.8\baselineskip}
\addtolength{\voffset}{-1.5\baselineskip}

\encabezado

%%%%%%%%%%%%%%%%%%%%%%%%%%%%%%%%%%%%%%%%
\section{Resultados de aprendizaje} 
%%%%%%%%%%%%%%%%%%%%%%%%%%%%%%%%%%%%%%%%

\begin{itemize}[leftmargin=*]
\item 
    \textbf{RdA 1:} Comprender los conceptos
fundamentales del Álgebra Lineal, incluyendo el estudio de matrices, determinantes, sistemas de ecuaciones lineales y espacios vectoriales, destacando su importancia en el análisis y resolución de problemas matemáticos.
    \begin{itemize}[leftmargin=*]
        \item \textbf{Criterio 1.1:}  Identifica las propiedades fundamentales de las matrices, incluyendo tipos, operaciones y la inversión de matrices.
        \item \textbf{Criterio 1.2:} Identifica los elementos de los sistemas de ecuaciones lineales y la clasificación de estos por su tipo de soluciones.
        \item \textbf{Criterio 1.3:} Comprende los conceptos de espacios vectoriales y aplicaciones lineales, incluyendo la base y dimensión, transformaciones lineales y sus propiedades.
    \end{itemize}
\item 
    \textbf{RdA 2:} Resolver problemas prácticos y teóricos relacionados con el Álgebra Lineal, utilizando técnicas de cálculo con matrices, determinantes y sistemas de ecuaciones, así como en el análisis de espacios vectoriales.
    \begin{itemize}[leftmargin=*]
        \item \textbf{Criterio 2.1:} Resuelve operaciones con matrices, incluyendo productos y cálculo de inversas; además del cálculo determinantes de matrices y su significado en el contexto del Álgebra Lineal.
        \item \textbf{Criterio 2.2:} Aplica diferentes métodos para resolver sistemas de ecuaciones lineales, como el método de eliminación de Gauss y técnicas matriciales.
        \item \textbf{Criterio 2.3:} Resuelve operaciones en espacios vectoriales y aplicaciones lineales, incluyendo el cálculo de bases, núcleo, imagen, así como valores y vectores propios.
    \end{itemize}
\item
    \textbf{RdA 3:} Aplicar los principios y métodos del Álgebra Lineal en una variedad de contextos demostrando comprensión de su relevancia y utilidad en situaciones prácticas y teóricas.
    \begin{itemize}[leftmargin=*]
        \item \textbf{Criterio 3.1:} Aplica métodos de Álgebra Lineal para resolver sistemas de ecuaciones en contextos variados, demostrando habilidad para seleccionar y emplear el método más adecuado.
        \item \textbf{Criterio 3.2:} Aplica los conceptos de espacios vectoriales y transformaciones lineales para modelar y resolver problemas de diversos contextos.
    \end{itemize}
\end{itemize}

%%%%%%%%%%%%%%%%%%%%%%%%%%%%%%%%%%%%%%%%
\section{Contenidos generales} 
%%%%%%%%%%%%%%%%%%%%%%%%%%%%%%%%%%%%%%%%

\begin{itemize}
\item 
    Matrices
\item 
    Sistemas de ecuaciones
\item 
    Determinantes
\item 
    Espacios vectoriales
\item 
    Aplicaciones lineales
\item 
    Valores y vectores propios
\end{itemize}

%%%%%%%%%%%%%%%%%%%%%%%%%%%%%%%%%%%%%%%%
\section{Actividades de evaluación} 
%%%%%%%%%%%%%%%%%%%%%%%%%%%%%%%%%%%%%%%%

\begin{center}
\small
\begin{tabular}{C{1.4cm}L{5.5cm}L{8cm}}
    \toprule
    \textbf{Criterio} & \textbf{Actividad} & \textbf{Descripción} \\
    \midrule
    1.1
        & Cuestionario en línea 1 (100\%)
        & Evalúa la identificación de tipos de matrices, sus operaciones básicas y el proceso de inversión de matrices. \\ 
    \midrule
    1.2
        & Cuestionario en línea 2 (60\%)
        & Verifica la identificación de sistemas de ecuaciones lineales y su clasificación según el tipo de solución. \\\cmidrule{2-3}
        & Video-exposición 1 (40\%)
        & Presenta y explica la identificación de sistemas de ecuaciones lineales y su clasificación según el tipo de solución. \\
    \midrule
    1.3
        & Cuestionario en línea 3 (60\%)
        & Evalúa la comprensión de espacios vectoriales, base, dimensión y propiedades de aplicaciones lineales. \\\cmidrule{2-3}
        & Video-exposición 2 (40\%)
        & Expone conceptos de transformaciones lineales y su relación con base y dimensión. \\
    \midrule
    2.1
        & Examen 1 - parte 1 (100\%)
        & Mide la resolución de operaciones con matrices, cálculo de inversas y determinantes en problemas aplicados. \\
    \midrule
    2.2
        & Examen 1 - parte 2 (100\%)
        & Evalúa la aplicación de métodos como eliminación de Gauss y técnicas matriciales para resolver sistemas lineales. \\
    \midrule
    2.3
        & Examen 2 (100\%)
        & Valora la resolución de problemas de espacios vectoriales, núcleo, imagen, valores y vectores propios. \\
    \midrule
    3.1
        & Reto 1 (100\%)
        & Aplica métodos de álgebra lineal para modelar y resolver sistemas en contextos prácticos diversos. \\
    \midrule
    3.2
        & Reto 2 (100\%)
        & Integra espacios vectoriales y aplicaciones lineales para modelar situaciones de diferentes contextos. \\
    \bottomrule
\end{tabular}
\end{center}



\begin{landscape}
%%%%%%%%%%%%%%%%%%%%%%%%%%%%%%%%%%%%%%%%
\section{Cronograma de Desarrollo del Curso} 
%%%%%%%%%%%%%%%%%%%%%%%%%%%%%%%%%%%%%%%%

\begin{center}\small
\setlength{\extrarowheight}{0ex}
\setlength{\belowrulesep}{.6ex}
\begin{longtable}{cccL{12cm}L{7.5cm}}
    \toprule
    &&\thead{Fecha}&\thead{Detalle de contenido} & \thead{Observación} \\
    \midrule
  \endfirsthead
    \multicolumn{5}{l}{\footnotesize \ldots viene de la página anterior}\\
    \toprule
    &&\thead{Fecha}&\thead{Detalle de contenido} & \thead{Observación} \\
    \midrule
  \endhead
        \bottomrule  \multicolumn{5}{r}{\footnotesize Continúa en la siguiente página\ldots}
  \endfoot
        \bottomrule
  \endlastfoot
1	&	1	&	16-mar	&	Matrices y operaciones con matrices	&		\\	
	&	2	&	18-mar	&	Sistemas de ecuaciones	&		\\	
	&	3	&	20-mar	&	Sistemas de ecuaciones	&		\\ \midrule	
2	&	4	&	23-mar	&	Inversa de matrices	&	Envío Reto 1	\\	
	&	5	&	25-mar	&	Determinantes: definición y propiedades	&		\\	
	&	6	&	27-mar	&	Desarrollo por cofactores, Aplicaciones	&		\\ \midrule	\rowcolor{celeste!50}
3	&	7	&	30-mar	&	Actividad de evaluación	&	Examen escrito 1	\\	
	&	8	&	01-abr	&	Vectores en el plano y en Rn	&	Cuestionario en línea 1	\\	\rowcolor{verde!50}
	&		&	03-abr	&		&	Feriado	\\ \midrule	
4	&	9	&	06-abr	&	Espacios vectoriales, Subespacios vectoriales	& Entrega Reto 1		\\	
	&	10	&	08-abr	&	Independencia linael, Conjuntos generadores	&	Entrega Video-Exp 1	\\	
	&	11	&	10-abr	&	Bases, Dimensión	&		\\ \midrule	
5	&	12	&	13-abr	&	Coordenadas y cambio de base	&		\\	
	&	13	&	15-abr	&	Suma de espacios, suma directa de espacios	&		\\	
	&	14	&	17-abr	&	Espacios con producto interno	&		\\ \midrule	
6	&	15	&	20-abr	&	Bases ortonormales, Complemento ortogonal	& Envío Reto 2		\\	
	&	16	&	22-abr	&	Aplicaciones lineales, núcleo e imagen	&		\\	
	&	17	&	24-abr	&	Propiedades de aplicaciones lineales, Isomorfismos	&	Entrega Video-Exp 2	\\ \midrule	
7	&	18	&	27-abr	&	Matriz asociada a una aplicación lineal	&	Cuestionario en línea 2	\\	
	&	19	&	29-abr	&	Valores y vectores propios de una matriz	&		\\	\rowcolor{verde!50}
	&		&	01-may	&		&	Feriado	\\ \midrule	
8	&	20	&	04-may	&	Valores y vectores propios de una aplicación lineal	&	Entrega Reto 2	\\	
	&	21	&	06-may	&	Diagonalización de matrices	&		\\	\rowcolor{celeste!50}
	&	22	&	08-may	&	Actividad de evaluación	&	Examen escrito 2	\\ 
\end{longtable}
\end{center}
\end{landscape}

\end{document} 